\documentclass[../main.tex]{subfiles}
\begin{document}

Chương này trình bày toàn bộ quá trình phân tích yêu cầu và thiết kế hệ thống triage log tự động. Xuất phát từ các cơ sở lý thuyết được đặt nền ở Chương 2, mục tiêu của chương này là chuyển hóa các mô hình toán học thành một kiến trúc phần mềm cụ thể, có thể triển khai và đánh giá được. Với mỗi lớp xử lý, không chỉ thiết kế luồng dữ liệu mà còn phải xác định các giá trị tham số cụ thể dựa trên cơ sở lý thuyết và phân tích thực nghiệm, thay vì để chúng như các ``magic number'' không có căn cứ.

\section{Phân tích Yêu cầu Hệ thống}
\label{section:3.1}

\subsection{Yêu cầu Chức năng}
\label{subsection:3.1.1}

Hệ thống được thiết kế để giải quyết bài toán triage tự động trên log Sysmon Windows theo một luồng xử lý hoàn toàn offline. Yêu cầu chức năng cốt lõi bao gồm bốn nhóm chính. Thứ nhất, hệ thống phải nhận vào một thư mục chứa các file \texttt{.evtx} (có thể từ nhiều máy hoặc nhiều ngày), sau đó xuất ra cây tiến trình có phân tầng rủi ro đầy đủ. Thứ hai, mỗi node trong cây phải được gắn điểm số phân rã thành sáu thành phần P, C, S, Seq, Col, CI — tức là analyst đọc output có thể truy xuất ngay lý do một tiến trình bị flag mà không cần tra cứu thêm. Thứ ba, hệ thống phải hỗ trợ log đa file một cách minh bạch: các file từ nhiều máy hoặc nhiều phiên ghi log được tự động merge theo timeline UTC, đảm bảo tính toàn vẹn của chuỗi cha-con xuyên file. Cuối cùng, mọi Sigma rule match phải được gắn nhãn MITRE ATT\&CK tương ứng trong output, giúp analyst ánh xạ ngay sang kill chain.

\subsection{Yêu cầu Phi chức năng}
\label{subsection:3.1.2}

Yêu cầu phi chức năng đặt ra bốn ràng buộc cứng phản ánh môi trường triển khai thực tế. Về hiệu năng, hệ thống phải xử lý file \texttt{.evtx} kích thước $\ge 1$ GB trong thời gian tính bằng phút — yêu cầu này loại bỏ ngay các cách tiếp cận parse XML thuần Python vốn là bottleneck chính. Về khả năng giải thích (XAI), mọi \texttt{Final\_Score} phải phân rã được thành các thành phần toán học tường minh; mô hình black-box không được chấp nhận vì analyst cần đủ cơ sở để đưa kết quả vào báo cáo điều tra. Về tính tự học, hệ thống không được yêu cầu bất kỳ dữ liệu baseline bên ngoài: toàn bộ profile ``bình thường'' phải được học từ chính log đang phân tích (\textit{on-the-fly self-learning}). Cuối cùng, về toàn vẹn pháp y, không một node nào mang thông tin hành vi hoặc novelty được phép bị xóa trong bất kỳ bước pruning nào — đây là ràng buộc trực tiếp từ nguyên lý Lossless Structural Deduplication.

\section{Tổng quan Kiến trúc}
\label{section:3.2}

\subsection{Triết lý Funnel Architecture}
\label{subsection:3.2.1}

Thách thức kỹ thuật cốt lõi khi xây dựng hệ thống phân tích log là xung đột giữa hai ràng buộc trái chiều: độ chính xác đòi hỏi thuật toán phức tạp chạy trên toàn bộ dữ liệu, trong khi hiệu năng đòi hỏi giảm thiểu số lượng phép tính. Cách giải quyết là áp dụng triết lý \textit{Funnel Architecture} — xếp các lớp xử lý theo thứ tự chi phí tính toán tăng dần, mỗi lớp chỉ truyền xuống lớp sau tập dữ liệu đã được tinh lọc. Lớp L0 và L1 thực hiện lọc thô với chi phí $O(N)$ — loại bỏ event không thuộc 9 EID cần thiết và quét Sigma rule. Lớp L2 xây dựng cấu trúc đồ thị và baseline với chi phí $O(N \log N)$, đồng thời nén nhiễu qua Lossless Structural Deduplication. Lớp L3 chạy các thuật toán nặng nhất (HDBSCAN, Markov scoring, MAD normalization) nhưng chỉ trên tập node đã được lọc, số lượng nhỏ hơn nhiều so với N sự kiện ban đầu. Lớp L4 tổng hợp kết quả và render output.

% Hình 3.1 — Kiến trúc pipeline 5 tầng
\begin{figure}[htbp]
\centering
\includegraphics[width=0.95\textwidth]{Hinhve/2.1.png}
\caption{Kiến trúc pipeline 5 tầng của hệ thống triage tự động}
\label{fig:architecture}
\end{figure}

Triết lý này cũng quyết định luồng dữ liệu một chiều: L0 $\to$ L1 $\to$ L2 $\to$ L3 $\to$ L4. Không có vòng phản hồi (feedback loop), không có lớp nào đọc kết quả của lớp sau. Điều này không chỉ đơn giản hóa thiết kế mà còn đảm bảo tính deterministic: với cùng một input, hệ thống luôn cho cùng một output.


\subsection{Quyết định Công nghệ: Rust (L0) + Python (L1–L4)}
\label{subsection:3.2.2}

Quyết định sử dụng hai ngôn ngữ xuất phát từ phân tích bottleneck thực nghiệm. Khi thử nghiệm với pipeline parse \texttt{.evtx} thuần Python sử dụng thư viện \texttt{python-evtx}, khoảng 70\% tổng wall-clock time nằm ở bước đọc và parse file nhị phân — không phải ở các thuật toán phân tích. Nguyên nhân: Python với GIL (Global Interpreter Lock) không thể tận dụng đa luồng thực sự cho I/O binding task; đồng thời, parse XML trung gian từ EVTX record phát sinh chi phí Serialize/Deserialize không cần thiết.

\textbf{Rust} giải quyết vấn đề này theo hai hướng. Thứ nhất, Rust có thể tận dụng đa luồng thực sự để đọc song song nhiều file \texttt{.evtx}. Thứ hai, thư viện \texttt{evtx} viết bằng Rust parse trực tiếp binary template mà không đi qua tầng XML trung gian.

Giao tiếp giữa Rust và Python được thực hiện qua \textbf{PyO3} với cơ chế \textit{zero-copy memory transfer}: thay vì Rust xuất kết quả ra JSON qua stdout rồi Python đọc vào (Serialize $\to$ IPC $\to$ Deserialize — ba lần copy), Rust tạo ra Python Dictionary object trực tiếp trong Rust heap, sau đó chuyển ownership sang Python heap thông qua PyO3 ABI. Python nhận được native dict object mà không cần bất kỳ bước parse nào. Cách tiếp cận này loại bỏ hoàn toàn bottleneck I/O giữa L0 và L1.

Python (L1–L4) phù hợp cho các lớp còn lại vì: hệ sinh thái thư viện khoa học dữ liệu (NumPy, scikit-learn, hdbscan) trưởng thành và được tối ưu tốt cho vectorized computation; việc viết và điều chỉnh các thuật toán thống kê phức tạp dễ dàng hơn nhiều so với Rust; và bottleneck đã được giải quyết ở L0 nên L1–L4 không cần hiệu năng của compiled language.

% Hình 3.2 — DFD Level 0: System boundary
\begin{figure}[htbp]
\centering
\includegraphics[width=0.5\textwidth]{Hinhve/2.2.png}
\caption{Data Flow Diagram mức 0 — Ranh giới hệ thống và luồng dữ liệu tổng thể}
\label{fig:dfd0}
\end{figure}

\section{Thiết kế Layer 0 — Data Ingestion \& Parser}
\label{section:3.3}

% Hình 3.3 — Sequence diagram L0
\begin{figure}[htbp]
\centering
\includegraphics[width=1.0\textwidth]{Hinhve/2.3.png}
\caption{Sequence diagram Layer 0: từ thư mục \texttt{.evtx} đến luồng event Python}
\label{fig:l0_seq}
\end{figure}

\subsection{Chiến lược Multi-file Merge}
\label{subsection:3.3.1}

Khi đầu vào là một thư mục chứa nhiều file \texttt{.evtx}, L0 thực hiện ba bước theo thứ tự. Bước đầu tiên là quét thư mục để liệt kê toàn bộ file \texttt{.evtx}, sau đó đọc metadata timestamp của từng file (timestamp sự kiện đầu và cuối) mà không parse toàn bộ nội dung. Bước này có chi phí $O(M)$ với $M$ là số file. Bước thứ hai là sắp xếp các file theo timestamp và xác định thứ tự parse để đảm bảo luồng event đầu ra có thứ tự UTC toàn cục. Bước thứ ba là parse song song các file và merge-sort luồng event theo UTC timestamp.

Lý do thứ tự UTC toàn cục là bắt buộc: mô hình Markov bậc 2 ở L3 xây dựng transition matrix dựa trên thứ tự cha-con trong cây tiến trình. Nếu tiến trình cha được ghi trong file \texttt{host1\_day1.evtx} và tiến trình con được ghi trong \texttt{host1\_day2.evtx}, mà hai file này không được merge đúng thứ tự, cạnh cha-con sẽ bị đảo ngược hoặc mất — khiến transition matrix sai về cơ bản và Seq score trở nên vô nghĩa.

\subsection{Path Normalization — Hai chiều ngược nhau}
\label{subsection:3.3.2}

Log Windows có thể biểu diễn cùng một đường dẫn theo nhiều cách khác nhau: \texttt{C:\textbackslash Windows\textbackslash System32\textbackslash cmd.exe}, \texttt{c:\textbackslash windows\textbackslash system32\textbackslash CMD.EXE}, và \texttt{\%SystemRoot\%\textbackslash System32\textbackslash cmd.exe} đều chỉ cùng một binary. Nếu không chuẩn hóa, hệ thống tạo ra ba node khác nhau trong đồ thị, làm sai lệch toàn bộ thống kê tần suất và phân cụm.

L0 áp dụng chuẩn hóa \textit{chiều đi}: chuyển toàn bộ path về lowercase, thay dấu \texttt{\textbackslash} bằng \texttt{/}, và giải nén biến môi trường (\texttt{\%SystemRoot\%} $\to$ \texttt{c:/windows}). Kết quả: \texttt{c:/windows/system32/cmd.exe} là biểu diễn duy nhất cho mọi variant.

Tuy nhiên, khi L1 thực hiện Sigma matching, xuất hiện vấn đề theo chiều ngược: hầu hết quy tắc SigmaHQ được viết theo convention Windows với backslash (\texttt{c:\textbackslash windows\textbackslash...}). Nếu so khớp trực tiếp chuỗi đã chuẩn hóa với rule dùng backslash, toàn bộ matching thất bại. Do đó, L1 phải thực hiện chuẩn hóa \textit{chiều về} — convert trở lại về backslash format trước khi so khớp Sigma, rồi lưu kết quả matching vào node theo định danh chuẩn hóa.

\section{Thiết kế Layer 1 — Sigma Engine}
\label{section:3.4}

% Hình 3.4 — Flow diagram L1
% \begin{figure}[H]
% \centering
% % \includegraphics[width=0.8\textwidth]{Hinhve/2.4.png}
% \caption{Luồng xử lý Layer 1: quét Sigma rule và tổng hợp điểm S}
% \label{fig:l1_flow}
% \end{figure}

\subsection{Thiết kế Confidence Override System}
\label{subsection:3.4.1}

Như đã phân tích ở Chương 2 (§2.4.6), điểm S được tính theo công thức Max Override với hệ số tin cậy $CM_i$. Giá trị mặc định của $CM_i$ được thiết lập theo severity level của rule, phản ánh thực tế rằng rule có severity thấp hơn thường có tỷ lệ False Positive cao hơn trong môi trường production:

\begin{table}[h]
\centering
\begin{tabular}{|l|c|c|}
\hline
\textbf{Severity} & \textbf{Base score $B_i$} & \textbf{$CM_i$ mặc định} \\ \hline
Critical & 8.0 & 1.00 \\ \hline
High     & 6.0 & 0.85 \\ \hline
Medium   & 4.0 & 0.60 \\ \hline
Low      & 2.0 & 0.40 \\ \hline
\end{tabular}
\caption{Giá trị mặc định của base score và Confidence Multiplier theo severity}
\label{tab:sigma_defaults}
\end{table}

Bên cạnh giá trị mặc định, hệ thống hỗ trợ ghi đè per-rule thông qua file \texttt{confidence\_override.json} ánh xạ rule ID sang $CM_i$ tùy chỉnh. Cơ chế này cho phép analyst điều chỉnh độ nhạy của từng rule mà không cần sửa rule gốc — ví dụ, rule ``Execution Of Non-Existing File'' có thể được hạ $CM_i$ xuống 0.3 trong môi trường hay sinh FP do phần mềm bảo mật, trong khi vẫn giữ nguyên logic phát hiện.

\subsection{Chọn $\lambda_S = 0.25$ — Logarithmic Corroboration Bonus}
\label{subsection:3.4.2}

Tham số $\lambda_S$ trong công thức $S = \min(10, \max_i(B_i \cdot CM_i) + \lambda_S \cdot \log_2(1 + N_{rules}))$ kiểm soát mức độ tác động của việc có nhiều rule khớp. Cần chọn $\lambda_S$ sao cho: (1) bonus có ý nghĩa — nhiều rule khớp tốt hơn một rule; (2) bonus không làm đảo lộn thứ tự severity — nhiều rule Low không được vượt một rule Critical.

Phân tích trường hợp biên: với Critical rule ($B = 8.0, CM = 1.0$), điểm nền = 8.0. Để 10 rule Low ($B = 2.0, CM = 0.5$) vượt được, cần $10 \times 0.5 \cdot 1.0 + \lambda_S \cdot \log_2(11) > 8.0$, tức $5 + 3.46\lambda_S > 8.0$, tức $\lambda_S > 0.87$. Chọn $\lambda_S = 0.25 \ll 0.87$ đảm bảo rule Critical luôn dominant. Với 10 rule Low, bonus = $0.25 \times \log_2(11) \approx 0.86$ — có ý nghĩa nhưng không phá vỡ thứ tự.

\section{Thiết kế Layer 2-A — Graph Build \& Deduplication}
\label{section:3.5}

% Hình 3.5 — Graph construction and deduplication
\begin{figure}[htbp]
\centering
\includegraphics[width=1.0\textwidth]{Hinhve/2.5.png}
\caption{Luồng dữ liệu Layer 2-A: Xây dựng đồ thị tiến trình và nén dữ liệu cấu trúc}
\label{fig:l2a_graph}
\end{figure}

Lớp 2-A chịu trách nhiệm tái tạo đồ thị tiến trình có hướng (DAG) từ luồng sự kiện EID 1 (Process Create) và EID 5 (Process Terminate), sử dụng \texttt{ProcessGuid} làm khóa chính thay vì \texttt{ProcessId} (PID) do đặc tính dễ bị hệ điều hành tái sử dụng của PID. Sau khi đồ thị thô được xây dựng, hệ thống áp dụng thuật toán Lossless Structural Deduplication để giảm kích thước đồ thị (tối ưu hiệu năng cho L3) mà vẫn bảo toàn các bằng chứng pháp y quan trọng.

\subsection{Chọn $n_{rep} = 2$ — Cân bằng giữa Nén và Độ phân giải Hành vi}
\label{subsection:3.5.1}

Thuật toán Deduplication gộp các tiến trình có cùng tên (image) chạy dưới cùng một tiến trình cha (cùng \texttt{ParentProcessGuid}) và có cùng trạng thái tín hiệu (không có Sigma match, không có hành vi mạng/file/registry). Tham số $n_{rep}$ quy định số lượng tiến trình đại diện được giữ lại cho mỗi nhóm tương đương. 

Việc chọn $n_{rep}$ mang tính đánh đổi (trade-off) quan trọng. Nếu $n_{rep} = 1$, đồ thị được nén tối đa, nhưng hệ thống sẽ đánh mất thông tin về phương sai nội bộ nhóm, làm cho thành phần IPAS (chấm điểm bất thường nội bộ) kém chính xác do tập mẫu quá nhỏ. Nếu $n_{rep} = 3$ hoặc lớn hơn, lợi ích nén giảm đáng kể, khiến L3 phải xử lý một lượng lớn các tiến trình tương đương, gây lãng phí tài nguyên. Thực nghiệm đo lường phần trăm giảm số lượng node và sai số IPAS trên các tập dữ liệu mẫu cho thấy $n_{rep} = 2$ mang lại điểm cân bằng tối ưu: giữ được đủ mẫu để đánh giá phương sai thống kê trong khi vẫn giảm được lượng lớn dữ liệu nhiễu lặp (ví dụ: \texttt{svchost.exe} sinh ra hàng loạt).

\subsection{Danh sách SHELL\_WRAPPERS}
\label{subsection:3.5.2}

Để ngăn chặn thuật toán vô tình gộp các tiến trình vỏ (shell) hoặc ngôn ngữ kịch bản có thể ẩn giấu payload nguy hiểm trong command line, hệ thống duy trì một danh sách các \texttt{SHELL\_WRAPPERS} được ưu tiên giữ lại (bỏ qua bước gộp). Danh sách này bao gồm: \texttt{cmd.exe}, \texttt{powershell.exe}, \texttt{pwsh.exe}, \texttt{bash.exe}, và \texttt{wsl.exe}. Cần phân biệt rõ: đây là các chương trình shell thực thi kịch bản độc đoán, khác với các công cụ quản trị như \texttt{wmic.exe} hay \texttt{certutil.exe}. Các công cụ quản trị (mặc dù bị attacker lạm dụng) vẫn tuân theo quy tắc nén thông thường nếu chúng không thực hiện hành vi đáng ngờ (như tải file qua mạng) hoặc không mang tính duy nhất về cấu trúc (Structural Novelty).

\section{Thiết kế Layer 2-B — Baseline Builder}
\label{section:3.6}

% Hình 3.6 — Baseline Builder flow
\begin{figure}[htbp]
\centering
\includegraphics[width=1.0\textwidth,height=0.4\textheight,keepaspectratio]{Hinhve/2.6.png}
\caption{Luồng dữ liệu Layer 2-B: Xây dựng baseline hành vi động}
\label{fig:l2b_baseline}
\end{figure}

Lớp 2-B chạy song song với 2-A, chịu trách nhiệm tự học baseline hành vi từ tập dữ liệu đang phân tích. Toàn bộ tiến trình được vector hóa thành không gian 10 chiều ($V_{10}$) và phân cụm bằng HDBSCAN.

\subsection{Tham số HDBSCAN: \texttt{min\_cluster\_size} và Chiến lược Sampling}
\label{subsection:3.6.1}

Thuật toán HDBSCAN yêu cầu một tham số cốt lõi: \texttt{min\_cluster\_size}. Thử nghiệm (\textit{sweep}) giá trị từ 5 đến 50 trên tập dữ liệu benchmark cho thấy: dưới 10, đồ thị phân mảnh thành quá nhiều vi-cụm (micro-clusters) vô nghĩa; trên 20, các nhóm hành vi nhỏ nhưng có ý nghĩa bị gộp chung vào cụm lớn, làm giảm độ nhạy (\textit{sensitivity}) trong phát hiện. Giá trị \texttt{min\_cluster\_size} = 15 mang lại điểm \textit{silhouette score} cao nhất.

Đối với các log dung lượng lớn, ma trận khoảng cách $N \times N$ của HDBSCAN có thể gây tràn RAM (Out-Of-Memory). Thử nghiệm thực tế đo lường bộ nhớ cho thấy HDBSCAN vượt mốc 8GB RAM tại ngưỡng $N \approx 60,000$. Do đó, hệ thống áp dụng ngưỡng lấy mẫu (\textit{sampling threshold}) là $50,000$ tiến trình (với 20\% buffer an toàn). Khi $N > 50,000$, hệ thống sẽ lấy mẫu phân tầng (\textit{stratified sample}) 50,000 tiến trình để huấn luyện HDBSCAN, sau đó dùng hàm \texttt{approximate\_predict()} để gán nhãn cụm cho các tiến trình còn lại.

Trong trường hợp dữ liệu phân hóa thành quá nhiều cụm (vượt quá mức kỳ vọng cho một môi trường Windows thông thường), hệ thống đặt một mức trần số cụm (Cluster Ceiling) $K_{max} = 30$. Nếu thuật toán trả về số cụm lớn hơn 30, các vi-cụm nhỏ nhất sẽ được hợp nhất vào cụm gần nhất dựa trên khoảng cách Euclidean giữa các centroid. Điều này giúp ổn định dữ liệu và tránh tình trạng nhiễu cụm.

\subsection{Thiết kế Trọng số Cụm Hành vi $\vec{w}_{sus}$}
\label{subsection:3.6.2}

Độ rủi ro của mỗi cụm ($risk_k$) được tính bằng tích vô hướng giữa centroid của cụm và vector trọng số nghi vấn $\vec{w}_{sus}$. Bảng \ref{tab:w_sus} liệt kê các trọng số được thiết kế dựa trên mức độ tương quan với các kỹ thuật tấn công trong MITRE ATT\&CK:

\begin{table}[h]
\centering
\begin{tabular}{|l|c|p{7cm}|}
\hline
\textbf{Đặc trưng} & \textbf{$w_{sus}$} & \textbf{Cơ sở lý luận} \\ \hline
$v_2$: \texttt{has\_injection} & 4.0 & Rất ít ứng dụng hợp lệ thực hiện Code Injection; tín hiệu cực mạnh của T1055. \\ \hline
$v_1$: \texttt{has\_network} & 3.0 & Phổ biến ở app hợp lệ nhưng là thành phần thiết yếu của C2/Exfiltration (T1071, T1048). \\ \hline
$v_3$: \texttt{has\_file\_drop} & 2.5 & Đặc trưng của Dropper và thiết lập Persistence. \\ \hline
$v_6$: \texttt{nif\_pair} & 2.0 & Mối quan hệ cha-con hiếm gặp là chỉ dấu mạnh của LOLBins. \\ \hline
$v_4$: \texttt{has\_reg\_mod} & 1.5 & Dùng cho Persistence (T1547) nhưng có nhiều use case hợp lệ (Cài đặt, Cập nhật). \\ \hline
$v_5$: \texttt{nif\_binary} & 1.5 & Binary ít xuất hiện có rủi ro, nhưng yếu hơn tính hiếm của cặp cha-con. \\ \hline
$v_{10}$: \texttt{parent\_sigma\_bin} & 1.0 & Kế thừa một phần ngữ cảnh nghi ngờ từ tiến trình cha. \\ \hline
$v_7 - v_9$: Các đặc trưng phụ & 0.5 & Entropy, độ dài command line, số lượng tiến trình con mang tính hỗ trợ. \\ \hline
\end{tabular}
\caption{Trọng số $\vec{w}_{sus}$ cho từng chiều đặc trưng hành vi}
\label{tab:w_sus}
\end{table}

\section{Thiết kế Layer 2-C — Temporal Aggregation}
\label{section:3.7}

Trong thực tế phân tích log, một hành vi đơn lẻ có thể không đủ để coi là bất thường, nhưng sự tích tụ của nhiều hành vi nghi ngờ trong một khoảng thời gian ngắn là chỉ dấu mạnh của một cuộc tấn công đang diễn ra. Lớp 2-C chịu trách nhiệm tổng hợp các tín hiệu theo thời gian cho từng tiến trình (\textit{Temporal Aggregation}).

Hệ thống sử dụng mô hình suy giảm theo hàm mũ (\textit{Exponential Decay}) để tính điểm tích lũy ($CS$ — \textit{Cumulative Score}). Khi một sự kiện nghi ngờ mới xảy ra tại thời điểm $t$, điểm tích lũy được cập nhật:
\begin{equation}
    CS_{new} = CS_{prev} \cdot e^{-\lambda \Delta t} + Event\_Score \tag{2.1}
\end{equation}
Trong đó $\Delta t$ là khoảng thời gian kể từ sự kiện nghi ngờ trước đó, và $\lambda$ là hệ số suy giảm. Thiết kế này đảm bảo rằng các hành vi diễn ra dồn dập sẽ cộng hưởng điểm nhanh chóng, trong khi các hành vi rời rạc sẽ bị suy giảm theo thời gian, giúp giảm tỷ lệ False Positive từ các hoạt động quản trị hệ thống thông thường.

\section{Thiết kế Layer 3 — Risk Evaluator}
\label{section:3.8}

Lớp 3 chịu trách nhiệm chấm điểm độc lập cho từng tiến trình trên năm khía cạnh: Điểm (P), Ngữ cảnh (C), Luật (S - đã tính ở L1), Chuỗi (Seq), và Tập thể (Col), trước khi tổng hợp bằng Corroboration Index (CI).

\subsection{Cấu hình Tham số cho P Score}
\label{subsection:3.7.1}

Điểm P được tổng hợp từ MAD ($P_{MAD}$) và IPAS. Hệ thống chọn ngưỡng \textbf{MAD Gate ($\theta$) = 2.5}. Về mặt lý thuyết thống kê, $Z_{robust} = 2.5$ tương đương với việc nằm ngoài khoảng tin cậy 98.8\% của phân phối chuẩn. Thực nghiệm xác nhận: $\theta = 2.0$ dẫn đến tỷ lệ False Positive cao; $\theta = 3.0$ bỏ sót các dị thường mức vừa; $\theta = 2.5$ mang lại điểm cân bằng F1 cao nhất. Ngưỡng này cũng được sử dụng nhất quán cho \textbf{Seq Score} để đảm bảo sự tương đương về mặt thống kê.

Về tỷ trọng kết hợp, hệ thống chọn \textbf{$\alpha = 0.6$} trong công thức $P = \alpha \cdot P_{MAD} + (1-\alpha) \cdot IPAS$. Thực nghiệm cho thấy IPAS có khả năng phát hiện thêm trường hợp \textit{weaponized binary} mà MAD bỏ sót. Hệ số $\alpha = 0.6$ ưu tiên tín hiệu toàn cục vững chắc hơn, trong khi vẫn dành 40\% sức nặng cho phát hiện bất thường nội bộ.

Hệ số \textbf{IPAS Scale ($\kappa$) = 2.0} được xác định bằng cách phân tích phân phối khoảng cách $dist_{raw}$ trên tập dữ liệu có ground truth. Đặc biệt, tiến trình thuộc cụm nhiễu \textbf{Cluster -1} của HDBSCAN được nhận một điểm cộng \textbf{Bonus ($\delta$) = 3.0}.

\subsection{Cấu hình Tham số cho C, Seq, và Col Score}
\label{subsection:3.7.2}

\textbf{Điểm Ngữ cảnh (C):} Áp dụng \textit{one-sided Z-score} với hệ số scale \textbf{$\gamma = 3.33$}. Đối với các tiến trình thuộc \textbf{Cluster -1} thụ động, điểm khởi điểm là $c_0 = 3.0$ cộng với một giá trị tối đa $c_1 = 2.0$ phụ thuộc vào độ hiếm (NIF). Tổng điểm cao nhất $c_0 + c_1 = 5.0$ ngăn ngừa việc gán nhãn rủi ro cao cho các tiến trình chỉ mang tính nhiễu cấu trúc.

\textbf{Điểm Chuỗi (Seq):} Sử dụng mô hình Markov bậc 2 với hệ số \textbf{Laplace Smoothing ($\varepsilon_L$) = 0.05}. Giá trị này giúp tránh lỗi $\log(0)$ khi gặp các chuỗi tấn công hoàn toàn mới chưa từng xuất hiện.

\textbf{Điểm Tập thể (Col):} Áp dụng phân tích Coefficient of Variation (CoV) để phát hiện \textit{beaconing}. Yêu cầu tối thiểu \textbf{$N_{min}^{Col} = 5$} timestamp.

\subsection{Corroboration Index (CI) và Tham số Tổng hợp}
\label{subsection:3.7.3}

CI khuếch đại rủi ro khi có nhiều thành phần vượt ngưỡng ý nghĩa ($\tau_x$). Các ngưỡng được chọn dựa trên thống kê: $\tau_P = 6.0$, $\tau_C = 6.0$, $\tau_S = 3.0$, $\tau_{Seq} = 5.0$, $\tau_{Col} = 6.0$.
\begin{equation}
    CI_{raw} = 1.0 + \Delta_{CI} \cdot (N_{active} - 2) \tag{2.2}
\end{equation}
\begin{equation}
    CI = \text{clip}(CI_{raw},\ 0.85,\ 1.5) \tag{2.3}
\end{equation}

Hệ số \textbf{$\rho_{floor} = 0.85$} được áp dụng để giải quyết triệt để "vùng chết toán học" (\textit{Dead Zone}). Với $\rho_{floor} = 0.85$, điểm duy trì được phần lớn trọng số cơ bản ngay cả khi tín hiệu đơn lẻ không quá mạnh.

Tham số khuếch đại \textbf{$\Delta_{CI} = 0.1$} được chọn để CI tăng từ từ theo số lượng tín hiệu. Mức trần \textbf{$\rho_{ceil} = 1.5$} giới hạn sự khuếch đại.

Đặc biệt, cơ chế \textbf{Screaming Sigma Override} sẽ bỏ qua trọng số nếu $S \ge \eta_\sigma = 7.0$ (tương đương ít nhất một rule Critical có độ tin cậy cao), đảm bảo cảnh báo Known-Threat không bị pha loãng bởi các thuật toán thống kê.

\textbf{Trọng số Thành phần (Weights):} Kết quả \textit{ablation study} dẫn đến cấu hình tỷ trọng: $w_S = 0.35$ (tác động Recall mạnh nhất), $w_P = 0.25$, $w_{Seq} = 0.20$, $w_C = 0.20$. Thành phần Col được cộng dồn độc lập ($+0.10 \cdot Col$).
\begin{equation}
    Final\_Score = (w_P P + w_C C + w_S S + w_{Seq} Seq) \times CI + 0.10 \cdot Col \tag{2.4}
\end{equation}

\section{Thiết kế Layer 4 — Aggregation \& Output}
\label{section:3.8}

Lớp cuối cùng thực hiện lan truyền rủi ro (Backpropagation) và định dạng kết quả báo cáo (Output).

\subsection{Lan truyền rủi ro (Guilt Backpropagation) và Decay Factor}
\label{subsection:3.8.1}

Quá trình quét đồ thị từ dưới lên (Bottom-Up) giúp tái tạo Kill Chain bằng cách cập nhật điểm của tiến trình cha dựa trên điểm của tiến trình con, sử dụng hệ số suy giảm \textbf{$\beta = 0.85$}. Thiết lập này đảm bảo chuyên gia nhìn thấy rõ 2-3 cấp lineage mang rủi ro trước khi tín hiệu tắt dần.

\subsection{Phân cấp Tier và Smart Collapse}
\label{subsection:3.8.2}

Dựa trên phân tích phân phối Final Score của các tập dữ liệu đa dạng, hệ thống thiết lập bốn cấp độ rủi ro (Tier):
\begin{itemize}
    \item \textbf{Tier 1 (CRITICAL):} Final Score $\ge 7.0$ — Mối đe dọa xác định, yêu cầu ứng phó ngay lập tức.
    \item \textbf{Tier 2 (HIGH):} Final Score $\ge 5.0$ — Hành vi dị thường nghiêm trọng, cần điều tra ưu tiên.
    \item \textbf{Tier 3 (MEDIUM):} Final Score $\ge 1.5$ — Hành vi nghi ngờ, cần giám sát thêm.
    \item \textbf{Tier 4 (LOW/SAFE):} Final Score $< 1.5$ — Hành vi bình thường hoặc nhiễu cấu trúc.
\end{itemize}

Để tối ưu hóa trải nghiệm của chuyên gia phân tích (tránh hiện tượng báo cáo dài hàng chục nghìn dòng, "trôi màn hình"), hệ thống áp dụng kỹ thuật \textbf{Smart Collapse}. Bất kỳ nhánh cây nào (từ Root xuống toàn bộ các lá) mà không có bất kỳ một nút nào vượt ngưỡng 1.5 (toàn bộ nhánh nằm ở Tier 4) sẽ bị thu gọn (collapsed) thành một dòng duy nhất liệt kê tóm tắt ở cuối báo cáo. Kiểm chứng trên tập dữ liệu benchmark xác nhận 0\% True Positive bị ẩn đi bởi cơ chế Smart Collapse này.
\section{Tổng kết chương}
\label{section:3.9}

Chương này đã hiện thực hóa nền tảng toán học từ Chương 2 thành một kiến trúc phần mềm 5 lớp. Bằng cách thiết lập cấu trúc pipeline \textit{Funnel Architecture} (kết hợp sức mạnh của Rust và Python) và xác định cụ thể các tham số cốt lõi (từ $n_{rep}=2$ trong nén đồ thị đến $\beta=0.85$ trong lan truyền rủi ro), thiết kế này không chỉ giải quyết bài toán hiệu năng phân tích khối lượng lớn log mà còn giải quyết được các giới hạn của hệ thống đánh giá tĩnh (Static Scoring) hay các hệ thống học máy hộp đen (Black-box ML). Các tham số không được chọn ngẫu nhiên mà đều có cơ sở từ quan sát thực nghiệm và nguyên lý thiết kế: loại bỏ ``vùng chết toán học'' và bảo tồn bằng chứng ngoại biên (novelty). Chương tiếp theo sẽ trình bày quá trình thực nghiệm hệ thống này trên các tập dữ liệu mô phỏng kịch bản tấn công thực tế và đánh giá độ hiệu quả dựa trên các chỉ số Recall, Precision, và F1.

\end{document}